\clearpage
\gdef\otChapNum{0}\pagestyle{otFront}
\phantomsection\label{defis-debut}\pdfbookmark[0]{Défis d'approfondissement}{bm-defis}
\otPageTitre{Défis d'approfondissement}{Chercher plus loin avec les outils de Première}
\fontsize{11}{16}\selectfont
Les 210 exercices des chapitres installent les bases et les méthodes. Les
14 défis suivants portent le total à \textbf{224 exercices corrigés}.
Ils prolongent chaque thème par une question qui demande de prendre du recul,
de choisir une représentation ou de relier plusieurs notions.\par\vspace{4mm}
\textbf{Avancé} désigne une modélisation ou une synthèse soutenue.
\textbf{Très exigeant} ajoute une discussion, une preuve ou plusieurs changements de point de vue.
Ces indications sont des repères pédagogiques, pas un classement officiel.\par\vspace{4mm}
Les défis sont facultatifs pour un premier passage. Ils utilisent les notions
de Première ; les étapes qui dépasseraient les outils du cours sont explicitement admises.
Les durées proposées sont indicatives.\par\vspace{5mm}
\editionEtape{Chercher avant de lire la solution}
Faire un dessin ou un tableau, tester un cas simple, puis écrire ce qui manque
pour conclure. En cas de blocage, lire seulement le premier indice, puis reprendre
la recherche. Les indications sont séparées des énoncés et des corrigés.
\begin{center}\renewcommand{\arraystretch}{1.25}
\begin{tabular}{@{}lll@{}}\toprule
Défis & Thèmes & Préparation conseillée\\\midrule
D01 à D02 & Algèbre et suites & Chapitres 1, 2, 12, 13\\
D03 à D06 & Analyse et trigonométrie & Chapitres 1 à 6, 8, 12\\
D07 à D08 & Géométrie & Chapitres 1, 7, 8\\
D09 à D11 & Probabilités et statistiques & Chapitres 1, 9 à 11\\
D12 à D14 & Algorithmes et synthèses & Plusieurs chapitres, selon le défi\\\bottomrule
\end{tabular}\end{center}
\editionRetenir{Une solution incomplète mais justifiée est un progrès.
Après la correction, refaire la question clé sans regarder, puis expliquer pourquoi la méthode fonctionne.}

\editionPage{defi-ex-1}{D01 · Deux racines dans un intervalle}{Recherche et rédaction}{Avancé · 35 à 50 min · Chapitres 1 et 13\hfill\hyperref[defi-cor-1-1]{Corrigé p.~\pageref*{defi-cor-1-1}}}
\begin{blocEdition}{defi-ex-1}
Pour $m\in\R$, on pose $P_m(x)=x^2-2mx+m$. On cherche les paramètres pour
lesquels $P_m$ possède \textbf{deux racines distinctes, toutes deux dans $]0;2[$}.
\begin{questions}
\item Calculer le discriminant. Pour quels $m$ existe-t-il deux racines distinctes ?
\item Montrer que deux racines dans $]0;2[$ imposent $P_m(0)>0$, $P_m(2)>0$
et que l'abscisse du sommet appartient à $]0;2[$.
\item Réciproquement, expliquer pourquoi ces trois conditions, ajoutées à un
discriminant strictement positif, suffisent. On peut s'aider de la parabole.
\item Déterminer tous les paramètres recherchés. Étudier les deux valeurs frontières.
\item Les seules conditions $P_m(0)>0$ et $P_m(2)>0$ suffisent-elles ? Donner un contre-exemple.
\end{questions}
La figure fournit des exemples pour conjecturer ; la réponse doit être démontrée.
\label{figedition-defi-ex-1-1}\figureGeo{fr-d01}
\end{blocEdition}

\editionPage{defi-ex-2}{D02 · Le seuil d’une moyenne}{Recherche et rédaction}{Très exigeant · 45 à 65 min · Chapitres 2 et 12\hfill\hyperref[defi-cor-2-1]{Corrigé p.~\pageref*{defi-cor-2-1}}}
\begin{blocEdition}{defi-ex-2}
On définit $u_0=0$ et $u_{n+1}=\frac45u_n+60$ pour $n\geq0$.
Pour $n\geq1$, on note $S_n=\sum_{k=0}^{n-1}u_k$ et $A_n=S_n/n$.
\begin{questions}
\item À l'aide d'une suite auxiliaire géométrique, montrer que
$u_n=300\bigl(1-(4/5)^n\bigr)$. Étudier ses variations.
\item Calculer $S_n$, puis montrer que
$A_n=300-\frac{1500}{n}\bigl(1-(4/5)^n\bigr)$.
\item Expliquer pourquoi $A_n<u_n$, puis établir
$A_{n+1}-A_n=(u_n-A_n)/(n+1)$. En déduire les variations de $A_n$.
\item Écrire une fonction Python qui renvoie le premier $n\geq1$ tel que $A_n\geq240$.
Préciser ce que représentent les variables à chaque passage de la boucle.
\item Le programme conduit à $n=25$. Vérifier le seuil en calculant $A_{24}$ et $A_{25}$.
Pourquoi ces deux seuls contrôles suffisent-ils après la question 3 ?
\end{questions}
La moyenne porte sur les termes $u_0$ à $u_{n-1}$ : elle ne comprend pas $u_n$.
\end{blocEdition}

\editionPage{defi-ex-3}{D03 · Trois tangentes par un point}{Recherche et rédaction}{Très exigeant · 55 à 75 min · Chapitres 3, 4 et 1\hfill\hyperref[defi-cor-3-1]{Corrigé p.~\pageref*{defi-cor-3-1}}}
\begin{blocEdition}{defi-ex-3}
On considère $f(x)=x^3-3x$. Pour $t\in[-1;2]$, on note $T_t$ la tangente
à sa courbe au point d'abscisse $t$. Pour $b\in\R$, on pose $P_b(1;b)$.
\begin{questions}
\item Montrer que $T_t$ a pour équation $y=(3t^2-3)x-2t^3$.
\item Montrer que $T_t$ passe par $P_b$ si et seulement si $g(t)=b$,
où $g(t)=-2t^3+3t^2-3$.
\item Étudier $g$ sur $[-1;2]$ et calculer ses valeurs aux bornes et aux points critiques.
\item Selon $b$, déterminer le nombre de tangentes $T_t$ qui passent par $P_b$.
Soigner les cas où $b$ est une valeur extrême ou une valeur frontière.
\item Pour $b=-5/2$, factoriser $4t^3-6t^2+1$ et déterminer les trois
abscisses de contact, puis les équations des trois tangentes.
\end{questions}
Seules les tangentes dont l'abscisse de contact est dans $[-1;2]$ sont comptées.
\label{figedition-defi-ex-3-1}\figureGeo{fr-d03}
\end{blocEdition}

\editionPage{defi-ex-4}{D04 · La fenêtre la plus lumineuse}{Recherche et rédaction}{Avancé · 35 à 50 min · Chapitres 1 et 4\hfill\hyperref[defi-cor-4-1]{Corrigé p.~\pageref*{defi-cor-4-1}}}
\begin{blocEdition}{defi-ex-4}
Une fenêtre est formée d'un rectangle de largeur $2r$ et de hauteur $h$, surmonté
d'un demi-disque de rayon $r$. La longueur totale de son \textbf{contour extérieur} est de 20 m.
On suppose $r>0$ et $h>0$ et on veut maximiser l'aire vitrée.
\begin{questions}
\item Exprimer le périmètre. Pourquoi le diamètre commun au rectangle et au demi-disque
ne doit-il pas être ajouté au contour extérieur ?
\item Exprimer $h$ en fonction de $r$ et préciser exactement le domaine admissible.
\item Montrer que l'aire est $A(r)=20r-(2+\pi/2)r^2$.
\item Déterminer les dimensions optimales et l'aire maximale, sous forme exacte.
Quelle relation simple lie $h$ et $r$ à l'optimum ?
\item Expliquer pourquoi le maximum calculé correspond à une vraie fenêtre et non à une forme dégénérée.
\end{questions}
\label{figedition-defi-ex-4-1}\figureGeo{fr-d04}
\end{blocEdition}

\editionPage{defi-ex-5}{D05 · Encadrer sans logarithme}{Recherche et rédaction}{Très exigeant · 45 à 60 min · Chapitres 5, 2 et 12\hfill\hyperref[defi-cor-5-1]{Corrigé p.~\pageref*{defi-cor-5-1}}}
\begin{blocEdition}{defi-ex-5}
On veut encadrer $e^{1/2}$ avec des puissances de nombres rationnels.
\begin{questions}
\item Étudier $F(x)=e^x-1-x$ sur $\R$ et montrer que $e^x\geq1+x$,
avec égalité seulement pour $x=0$.
\item Pour $0<t<1$, appliquer ce résultat à $-t$, puis justifier
$1+t<e^t<1/(1-t)$.
\item Pour tout entier $n\geq1$, en déduire
\[L_n=\left(1+\frac1{2n}\right)^n<e^{1/2}<
U_n=\left(1-\frac1{2n}\right)^{-n}.\]
\item Programmer le calcul de $L_n$ et $U_n$, puis donner des encadrements
pour $n=10$, $100$ et $1000$. Arrondir la borne inférieure vers le bas et la borne supérieure vers le haut.
\item Que prouve l'encadrement théorique ? Que ne suffit pas à prouver la seule observation des décimales calculées ?
\end{questions}
On utilise uniquement les propriétés de l'exponentielle ; le logarithme n'est pas nécessaire.
\end{blocEdition}

\editionPage{defi-ex-6}{D06 · Additionner sinus et cosinus}{Recherche et rédaction}{Très exigeant · 45 à 65 min · Chapitres 6, 1 et 8\hfill\hyperref[defi-cor-6-1]{Corrigé p.~\pageref*{defi-cor-6-1}}}
\begin{blocEdition}{defi-ex-6}
Pour $t\in[0;2\pi[$, on pose $x=\cos t$, $y=\sin t$. On cherche à résoudre
$\cos t+\sin t=k$, où $k$ est un paramètre réel.
\begin{questions}
\item Traduire le problème comme une intersection entre un cercle et une droite.
\item En utilisant $(x-y)^2\geq0$, montrer que $-\sqrt2\leq x+y\leq\sqrt2$.
\item Lorsque $x+y=k$, montrer que $xy=(k^2-1)/2$, puis que $x$ et $y$
sont les racines de $X^2-kX+(k^2-1)/2=0$.
\item Classer le nombre de solutions $t$ selon $k$. Donner les coordonnées des points correspondants.
\item Résoudre explicitement pour $k=0$, $k=1$ et $k=\sqrt2$.
\end{questions}
L'intervalle est ouvert en $2\pi$ pour ne pas compter deux fois le même point du cercle.
\label{figedition-defi-ex-6-1}\figureGeo{fr-d06}
\end{blocEdition}

\editionPage{defi-ex-7}{D07 · Un triangle sous deux contraintes}{Recherche et rédaction}{Très exigeant · 40 à 60 min · Chapitres 7 et 8\hfill\hyperref[defi-cor-7-1]{Corrigé p.~\pageref*{defi-cor-7-1}}}
\begin{blocEdition}{defi-ex-7}
Dans un repère orthonormé, $A(-2;0)$, $B(2;0)$ et $M(x;y)$ vérifient
\[\vect{MA}\cdot\vect{MB}=0,\qquad MA^2+2MB^2=k.\]
Le paramètre $k$ est réel.
\begin{questions}
\item Montrer que la première condition équivaut à $x^2+y^2=4$.
\item Sous cette condition, établir $MA^2=8+4x$ et $MB^2=8-4x$.
\item En déduire les coordonnées possibles de $M$ en fonction de $k$ et discuter leur nombre.
\item Pour quels $k$ obtient-on un \textbf{vrai triangle} $AMB$ rectangle en $M$ ?
Pourquoi les valeurs frontières doivent-elles être exclues ?
\item Pour quel $k$ le triangle est-il aussi isocèle en $M$ ? Donner ses deux sommets possibles.
\end{questions}
\label{figedition-defi-ex-7-1}\figureGeo{fr-d07}
\end{blocEdition}

\editionPage{defi-ex-8}{D08 · Une famille de cercles tangents}{Recherche et rédaction}{Très exigeant · 55 à 75 min · Chapitres 8, 7 et 1\hfill\hyperref[defi-cor-8-1]{Corrigé p.~\pageref*{defi-cor-8-1}}}
\begin{blocEdition}{defi-ex-8}
Pour $m\in\R$, on considère $\mathcal C_m$ d'équation
$x^2+y^2-4x-2my=0$ et la droite $d:x+y=6$.
\begin{questions}
\item Déterminer le centre et le rayon de $\mathcal C_m$. Montrer que tous
ces cercles passent par $A(0;0)$ et $B(4;0)$.
\item Substituer $y=6-x$ dans l'équation du cercle. Montrer que le discriminant
de l'équation obtenue est $4(m^2+8m-8)$.
\item Classer, selon $m$, le nombre de points communs à $d$ et $\mathcal C_m$.
Pour quels paramètres y a-t-il tangence ?
\item On note $\Omega_m(2;m)$. Calculer son projeté orthogonal $H_m$ sur $d$.
Vérifier, pour chacun des paramètres de tangence, que $H_m$ appartient au cercle.
\item Expliquer pourquoi la droite $d$ est alors perpendiculaire au rayon au point de contact.
\end{questions}
La figure de l'énoncé représente seulement les points fixes et la droite ; les cercles restent à déterminer.
\label{figedition-defi-ex-8-1}\figureGeo{fr-d08-ex}
\end{blocEdition}

\editionPage{defi-ex-9}{D09 · Choisir une règle de contrôle}{Recherche et rédaction}{Très exigeant · 55 à 75 min · Chapitres 9 et 10\hfill\hyperref[defi-cor-9-1]{Corrigé p.~\pageref*{defi-cor-9-1}}}
\begin{blocEdition}{defi-ex-9}
Dans un modèle fictif, $10\,\%$ des articles présentent un défaut $D$.
On applique deux tests $T_1$ et $T_2$, de résultat positif $+$ ou négatif $-$.
Sur les articles défectueux, chaque test est positif avec probabilité $0{,}9$ ;
sur les autres, avec probabilité $0{,}2$.
\textbf{Dans chacun de ces deux groupes}, le résultat du second test a les mêmes
probabilités quel que soit le résultat du premier.
\begin{questions}
\item Compléter l'arbre à trois niveaux. Les deux tests sont-ils pour autant
indépendants dans la population entière ? Comparer $P(T_1^+\cap T_2^+)$ à $P(T_1^+)P(T_2^+)$.
\item Calculer la probabilité du défaut sachant qu'un test est positif,
puis sachant que les deux sont positifs.
\item La règle U écarte un article dès qu'au moins un test est positif.
La règle B ne l'écarte que si les deux sont positifs. Pour un article choisi dans toute la production,
calculer, pour chaque règle, $P(\overline D\text{ et rejet})$ et $P(D\text{ et acceptation})$.
\item Le premier type d'erreur coûte 1 unité, le second $\lambda$ unités, avec $\lambda\geq0$.
Comparer les coûts moyens des deux règles selon $\lambda$.
\end{questions}
\label{figedition-defi-ex-9-1}\figureGeo{fr-d09}
\end{blocEdition}

\editionPage{defi-ex-10}{D10 · Même moyenne, objectifs différents}{Recherche et rédaction}{Très exigeant · 40 à 60 min · Chapitres 10 et 1\hfill\hyperref[defi-cor-10-1]{Corrigé p.~\pageref*{defi-cor-10-1}}}
\begin{blocEdition}{defi-ex-10}
Un jeu présente trois issues $I_1,I_2,I_3$ de probabilités $1/4,1/2,1/4$.
Pour un paramètre $a\in[0;2]$, le lot brut $L_a$ vaut respectivement
$2a$, $2-a$ et $6$ unités. Le prix d'entrée est $5/2$ ; le gain net est $G_a=L_a-5/2$.
\begin{questions}
\item Donner la loi issue par issue de $L_a$ et calculer son espérance.
À quelle valeur de $a$ deux lots coïncident-ils ? Comment regrouper alors les probabilités ?
\item Calculer $E(G_a)$ et montrer que
$V(G_a)=\frac32a^2-2a+\frac{19}{4}$.
\item Déterminer le paramètre qui minimise cette variance et la valeur minimale.
\item Calculer la probabilité $P(G_a>0)$ selon $a$, en traitant le seuil d'égalité.
Le choix de variance minimale maximise-t-il la probabilité de gagner strictement ?
\item Pour une prévision réelle $r$, développer $E((L_a-r)^2)$.
Déterminer simultanément les valeurs de $a$ et $r$ qui minimisent cette erreur quadratique moyenne.
\end{questions}
\label{figedition-defi-ex-10-1}\figureGeo{fr-d10}
\end{blocEdition}

\editionPage{defi-ex-11}{D11 · La dispersion d’un groupe réuni}{Recherche et rédaction}{Très exigeant · 45 à 65 min · Chapitres 11 et 1\hfill\hyperref[defi-cor-11-1]{Corrigé p.~\pageref*{defi-cor-11-1}}}
\begin{blocEdition}{defi-ex-11}
Le groupe A contient les valeurs 0 et 2 ; le groupe B contient 8 et 10.
La variance d'une série est la moyenne des carrés des écarts à sa moyenne.
\begin{questions}
\item Calculer les moyennes et variances des deux groupes, puis de leur réunion.
Pourquoi la variance de la réunion n'est-elle pas la moyenne des deux variances ?
\item Plus généralement, deux groupes ont des tailles $n_A,n_B>0$, des moyennes
$\mu_A,\mu_B$ et des variances $V_A,V_B$. Poser $p=n_A/(n_A+n_B)$ et exprimer
la moyenne $\mu$ de leur réunion.
\item Dans le groupe A, développer
$(x-\mu)^2=((x-\mu_A)+(\mu_A-\mu))^2$ puis moyenner.
Que devient le terme croisé ? Faire de même dans B.
\item En déduire la formule
\[V=pV_A+(1-p)V_B+p(1-p)(\mu_A-\mu_B)^2.\]
\item Deux groupes ont chacun une variance égale à 1. La variance de leur réunion
peut-elle être aussi grande que l'on veut ? Construire un exemple avec deux groupes de même taille.
\end{questions}
\label{figedition-defi-ex-11-1}\figureGeo{fr-d11}
\end{blocEdition}

\editionPage{defi-ex-12}{D12 · Un algorithme qui encadre une racine}{Recherche et rédaction}{Très exigeant · 45 à 60 min · Chapitres 12, 3 et 4\hfill\hyperref[defi-cor-12-1]{Corrigé p.~\pageref*{defi-cor-12-1}}}
\begin{blocEdition}{defi-ex-12}
On pose $f(x)=x^3+x-1$. \textbf{On admet l'existence} d'une solution $\alpha$ de
$f(x)=0$ dans $[0;1]$ ; il ne s'agit pas ici de démontrer un théorème de continuité.
\begin{questions}
\item Justifier que cette solution est unique.
\item On part de $a=0$, $b=1$, puis on calcule $m=(a+b)/2$.
Si $f(m)<0$, on remplace $a$ par $m$ ; sinon, on remplace $b$ par $m$.
Expliquer pourquoi $a\leq\alpha\leq b$ reste vrai après cette étape.
\item Écrire une fonction qui réalise exactement $N$ étapes et renvoie les deux bornes.
Préciser les entrées autorisées.
\item Montrer que la largeur finale vaut $2^{-N}$.
Quel est le plus petit $N$ qui garantit une largeur au plus $10^{-3}$ ?
\item Pour $N=10$, vérifier l'encadrement exact
$698/1024\leq\alpha\leq699/1024$ en étudiant les signes aux deux bornes.
Quelle erreur maximale garantit le milieu de cet intervalle ?
\end{questions}
\end{blocEdition}

\editionPage{defi-ex-13}{D13 · Une garantie pour tous les x}{Recherche et rédaction}{Très exigeant · 45 à 65 min · Chapitres 13, 1 et 4\hfill\hyperref[defi-cor-13-1]{Corrigé p.~\pageref*{defi-cor-13-1}}}
\begin{blocEdition}{defi-ex-13}
Pour $m\in\R$, on définit $P_m(x)=x^2-2mx+m$, et l'on s'intéresse à $x\in[0;1]$.
\begin{questions}
\item Formuler la négation de : « Pour tout $x\in[0;1]$, $P_m(x)\geq0$. »
\item Montrer que cette propriété impose $0\leq m\leq1$ en regardant les bornes.
Prouver la réciproque avec la forme canonique.
\item On note $q(m)$ le minimum de $P_m$ sur $[0;1]$. Déterminer $q(m)$
pour $m<0$, $0\leq m\leq1$ et $m>1$.
\item Quel paramètre fournit la plus grande garantie $q(m)$ ?
Donner une seconde preuve en utilisant la valeur commune $P_m(1/2)$.
\item Comparer les deux affirmations : « Il existe un $m$ tel que, pour tout $x$,
$P_m(x)\geq1/4$ » et « Pour tout $m$, il existe un $x$ tel que $P_m(x)\geq1/4$ ».
Préciser les domaines et justifier chacune des réponses.
\end{questions}
\end{blocEdition}

\editionPage{defi-ex-14}{D14 · Qualité, coût et risque}{Recherche et rédaction}{Très exigeant · 60 à 80 min · Chapitres 5, 9, 10 et 12\hfill\hyperref[defi-cor-14-1]{Corrigé p.~\pageref*{defi-cor-14-1}}}
\begin{blocEdition}{defi-ex-14}
Dans un atelier fictif, une intensité de contrôle $x\in[0;4]$ entraîne un coût
de $2x$ unités par article et une probabilité de défaut $p(x)=\frac15e^{-x}$.
Avant ce coût, un article rapporte 12 unités s'il est sans défaut et $-8$ sinon.
Le gain net après contrôle est noté $G_x$.
\begin{questions}
\item Donner les deux valeurs de $G_x$, leurs probabilités, puis montrer que
$E(G_x)=g(x)=12-2x-4e^{-x}$.
\item Montrer que le maximum est atteint pour l'unique $\alpha\in]0;1[$ tel que $e^\alpha=2$.
On admet l'existence de ce réel. Aucun logarithme n'est demandé.
\item Écrire un algorithme de dichotomie sur $[0;1]$ pour encadrer $\alpha$ à $10^{-3}$ près.
En déduire une approximation du gain moyen maximal.
\item Montrer que $V(G_x)=80e^{-x}-16e^{-2x}$ et étudier ses variations sur $[0;4]$.
Au point optimal pour l'espérance, calculer la probabilité de défaut et la variance exacte.
\item Deux articles sont contrôlés avec la même intensité $\alpha$ ; leurs défauts sont indépendants.
Dessiner l'arbre, puis calculer la probabilité d'au moins un défaut.
La plus grande intensité maximise-t-elle aussi le gain moyen ?
\end{questions}
\end{blocEdition}

